\section{\titleProofs}\label{sec:proofs}
  The previous sections were concerned with how cryptographers 
    can use the tools developed for the QROM to give
    post-quantum proofs of security. There, I did my best to leave the quantum
    analyses to surface-level only, taking the various lemmas as given, and focusing
    instead on developing an intuition for how to reason about the QROM.

  However, proving results about the QROM itself can serve as 
    rich playground for the more powerful tools found in \emph{Quantum
    Information Theory}. So if you, like
    many of us before you, have gotten a
    taste for all things quantum, and would
    like to dive deeper, then this is the section for you.

  In the following, I will introduce the powerful \emph{density matrix}
    formulation of quantum mechanics, introduced by legendary physicist and
    computer scientist John von Neumann, and show how, once understood, 
    it can provide elegant (and often simple!) proofs of the lemmas that we
    employed.

  \hhnote{
    In which I present all the tools needed to prove all the lemmas
    contained herein, some of which are then proven in-text and some of which are
    left as exercises (but with solutions).
  }
  
  \subsection{The Density Matrix Formalism}
    \hhnote{A primer.}

  \subsection{Tools of Quantum Information Theory}
    \hhnote{
      Distance measures, distinguishability, gentle measurements,
      almost-commuting; essentially a summary of Wilde chapter 9.
    }

  \subsection{Equivalent Oracles}
    \subsubsection{Superposition Oracles}
      \hhnote{
        A ``warmup'' proof showing that superposition oracles are
        equivalent to monolithic oracles.
      }

    \subsubsection{Superposition Oracles}
      \hhnote{
        The more involved proof of CROs
        being the same.
      }

  \subsection{One-way to Hiding}
    \subsubsection{The Original O2H Lemma}
      \hhnote{Include as warmup?}

    \subsubsection{The Semi-Classical O2H Lemma}
      \hhnote{
        Prove the most general version; essentially just repeat 
        the proof of the paper in your own notation.
      }
      
      Next, let's re-derive the Semi-Classical O2H bound by introducing
        $\measurement_\extev$ directly into the random oracle.
        If $\extractor$ has access to a $\qpca$ oracle, then they may 
        implement this oracle directly, which is shown in
        \cref{fig:semiclassical-oracle}.

      We again consider the setting that $\distinguisher$ is given some input $z$
        that may depend arbitrarily on $\realRO$, and that $\extractor$ provides a
        simulation for $\distinguisher$, except now it is implemented as a
        Semi-Classical CRO.

      \hhnote{
        The CRO is probably unnecessary; however, if we want $\RO_1$ to
        be related to $\RO_0$ in the way specified in the SC-O2H, then we
        could probably benefit from formulating it in terms of an independent
        oracle, and using the CSWAP technique.
        (Of course, the independent oracle \emph{could} then be implemented as
        a CRO.) \\
        Or wait \ldots the measurement projects the state onto the subspace
        independent of $\careset$ \ldots so we \emph{can} really just call
        $\RO_0$ instead?
      }

      \begin{figure}[t]
        \centering
        \begin{pchstack}[center]
          \procedure{Oracle $\idealRO(\inputreg, \outputreg : \bitreg)$}{
            \pclinecomment{1.\ Check input} \\
            \pcforeach \ket{x} \pcin\ \inputreg \pcdo \\
              \t \pcif x \in \careset \pcthen \applyU{\Xgate}{\bitreg} \\
            b \sample \meas{}{\bitreg} \\
            \pcif b = 1 \pcthen \pcabort \\
                %\t\t \target \sample \meas{}{\inputreg} \\
                %\t\t \pcabort \pcand \pcreturn \target \\
            \pclinecomment{2.\ Process query} \\
            (\inputreg, \outputreg) \gets \RO_0(\inputreg, \outputreg) \\
            \pcreturn (\inputreg, \outputreg)
          }
          
          %\pchspace

          %\procedure{Oracle $\RO[CRO](\inputreg, \outputreg : \database)$}{
          %  \applyU{\decomp}{\inputreg \concat \database} \\
          %  \applyU{\Hgate^{\otimes m}}{\outputreg} \\
          %  \applyU{\entangle}{\inputreg \concat \outputreg \concat \database} \\
          %  \applyU{\Hgate^{\otimes m}}{\outputreg} \\
          %  \applyU{\decomp}{\inputreg \concat \database} \\
          %  \pcreturn (\inputreg, \outputreg)
          %}
        \end{pchstack}
        \caption{
          \label{fig:semiclassical-oracle}
          The simulated semi-classical oracle $\idealRO$, implemented as a compressed
          oracle (recalled from \cref{def:compressed-ros}). Here, $\bitreg$ is a
          single-qubit register initialized to $\ket 0$, and $\database$ is
          the database, initialized to the list of $q$ identical tuples
          $(\emptystring, 0^m)$, (i.e.\ at the outset, the oracle runs the method
          $\database \gets \reginit{q}{\ket{\emptystring, 0^m}}$). 
        }
      \end{figure}

    \subsubsection{The Resampling Lemma}
      \hhnote{Include, or just rely on Joseph's permutations instead?}

    \subsubsection{The Adaptive Reprogramming Framework}
      \hhnote{Include, or just rely on Joseph's permutations instead?}

    \subsubsection{Joseph's Counterexample}
      \hhnote{
        Showing how the proof, which we saw took a lot of calculation to derive
        the state directly, becomes like two lines with the tools that we now
        have.
      }
